% !TeX program = pdflatex
% ALIUS Bulletin native reconstruction source.
% Compile from the ALIUS-bulletin project root. This file does not include a pre-existing article PDF.
\providecommand{\ALIUSRootPrefix}{}

\ifdefined\ALIUSIssueBuild
\else
\documentclass[a4paper,11pt]{article}
\usepackage[margin=22mm]{geometry}
\usepackage[T1]{fontenc}
\usepackage[utf8]{inputenc}
\usepackage[hidelinks]{hyperref}
\usepackage[protrusion=true,expansion=false]{microtype}
\usepackage{parskip}
\fi
\title{What is it like to be in early sensory life ?}
\author{An episode with Anna Ciaunica recorded by Daniel Friedman edited by Lena Coutrot, Daniel Friedman and Matthieu Koroma}
\date{ALIUS Bulletin 6}

\ifdefined\ALIUSIssueBuild
\else
\begin{document}
\fi
\maketitle
\section*{Metadata}
\textbf{Piece type:} podcast\par
\textbf{DOI:} 10.5281/zenodo.7457394\par
\textbf{Keywords:} Early Sensory Life, Consciousness, Multisensory, Embodiment, Shared Experience\par
\section*{Abstract}
In this episode, Anna Ciaunica reviews the recent progress on the scienfic study of conscious experiences in early sensory life.\par
\section*{Extracted Text}
\begingroup
\small
\noindent What is it like to be in early sensory life? An episode with Anna Ciaunica Recorded by Daniel Friedman and edited by Lena Coutrot and Mahieu Koroma Cite as: Ciaunica, A. (2023). What is it like to be in early sensory life ? ALIUS Bullen , 6 , hps://doi.org/10.5281/zenodo.7457394 Anna Ciaunica amaciaunica@fc.ul.pt Centro de Filosofia das Ciencias da Universidade de Lisboa, Portugal Abstract In this episode, Anna Ciaunica reviews the recent progress on the scienfic study of conscious experiences in early sensory life. keywords : Early Sensory Life, Consciousness, Mulsensory, Embodiment, Shared Experience The video and audio version of this episode are available on hps://aliusresearch.org/bullen06-ciaunica-podcast.html Consciousness is a mystery that many philosophers and scientists have attempted to tackle throughout the centuries. But what if instead of asking what consciousness is , one asks how conscious experiences emerge, develop and unfold in our life journey? Let's go back to 'suare one' and look at how experiences arise in early human life. What is it like to experience and perceive the world at the dawn of our lives? One basic yet overlooked aspect of current discussions in both philosophy and cognitive neuroscience is that experiences rst develop within another human body. e most basic perceptions and actions emerge already in utero . Crucially, while not all humans will have the experience of being pregnant or carrying a baby, the experience of perceiving and growing within another person's body is universal. Why is this important? How can this observation help us to understand conscious experiences? One basic yet overlooked aspect of current discussions in both philosophy and cognive neuroscience is that experiences first develop within another human body. ALIUS Bulletin n6 (2023) aliusresearch.org/bulletin Anna Ciaunica - What is it like to be in early sensory life ? 2 When I look outside my window, as an adult, I can see trees, hear the noise of the cars, smell the rain. All these perceptions form the content of my conscious experiences. But in the womb, starting from 22 weeks onwards, the world looks diherent. How? First, my entire body is immersed in a liuid environment, which makes my movements literally uid. My eyes are mostly closed, and even if the light comes in through my closed eyelids, my visual perception is limited. e proximal senses, such as touch and smell, are taking the lead in exploring the world. I can hear things, but most importantly, I can hear someone's heart beating all the time, continuously, rhythmically. e early sensory world of a human is not isolated nor solitary. It is a rich multisensory experiential orchestra with several players, involving others' embodied presence, experiences and moving bodies. Perceiving the world is not just about me, one individual looking at the world through a window. At the beginning of our lives, perceiving is about co-perceiving a shared world. We literally share time and space and our bodies with another human being. The early sensory world (...) is a rich mulsensory experienal orchestra with several players, involving others' embodied presence, experiences and moving bodies. is observation is in line with recent work in mind and brain research stipulating that our perceptions, cognitions and actions are fundamentally geared towards biological self-preservation, i.e., the need to maintain and regulate the psychological and physiological needs for the integrity of the living organism (the human body or individual) within a wider and highly volatile social and physical environment (Varela et al. 1991). is is something we share with all living organisms. In order to full the fundamental conditions for self-preservation, humans need to constantly move, act and interact with the physical and social environment. Now, one way the human organism may complete this task is described by the inuential 'Predictive Processing' framework (Friston 2010; Clark 2013; Hohwy 2013). e idea here is that we do not perceive the world as it is. Rather we perceive it as we 'expect' or 'predict' it will be, based on prior perceptions and experiences. ALIUS Bulletin n6 (2022) aliusresearch.org/bulletin Anna Ciaunica - What is it like to be in early sensory life ? 3 If this is so, then in order to understand the nature of conscious experience in the here and now of adult life, it is essential to look at how these experiences get oh the ground from the outset, in early life (Ciaunica et al. 2021a, 2021b). is is because adult conscious experiences cannot be addressed in isolation of prior (early life) perceptual experiences. Regardless of when various forms of consciousness rst emerge, there is an experiential continuum between early and later experiences. e key idea here is that by endorsing a bottom-up, sensory and developmental perspective in exploring how conscious experiences dynamically arise and develop in concert with the developing organism, we may gain important insights into what consciousness 'is' and its biological structure. Unveiling the mystery of consciousness perhaps then will depend less on what makes human experiences uniuely special, but rather on what connects them with life under all its forms. Adult conscious experiences cannot be addressed in isolaon of prior (early life) perceptual experiences. References Ciaunica, A., Constant, A., Preissl, H., Fotopoulo, K. (2021) e rst prior: from co-embodiment to co-homeostasis in early life. Consciousness and Cognition . 91:103117. https://doi.org/10.1016/j.concog.2021.103117 Ciaunica, A., Safron, A., \& Delaeld-Butt, J. (2021). Back to suare one: the bodily roots of conscious experiences in early life. Neuroscience of Consciousness , 2021 (2), niab037. https://doi.org/10.1093/nc/niab037 Clark, A. (2013) Whatever next? Predictive brains, situated agents, and the future of cognitive science. Behavioral and Brain Science . 36:181-204. https://doi.org/10.1017/S0140525X12000477 Friston, K. J. (2010) e free-energy principle: a unied brain theory? Nature Review Neuroscience . 11:127-38. https://doi.org/10.1038/nrn2787 Hohwy, J. (2013) e predictive mind . Oxford: Oxford University Press. Varela, F., ompson, E., Rosch E. (1991) e Embodied Mind: Cognitive Science and Human Experience . Cambridge, MA: MIT Press. ALIUS Bulletin n6 (2022) aliusresearch.org/bulletin\par
\endgroup

\ifdefined\ALIUSIssueBuild
\clearpage
\expandafter\endinput
\fi
\end{document}
